\chapter{Troponin}

\section*{What Is This Test?}
Cardiac troponin tests measure cardiac-specific isoforms of the troponin complex (troponin I or T) released into the circulation when cardiomyocytes are injured. Troponin is the preferred biomarker for myocardial injury, but an elevated result does not by itself diagnose myocardial infarction. MI requires a rise and/or fall in troponin with at least one value above the assay-specific 99th-percentile upper reference limit plus clinical evidence of ischemia. High-sensitivity assays support validated early rule-in/rule-out pathways, but these pathways and cutoffs are assay-specific.

\section*{Alternative Names}
Troponin I (cTnI); Troponin T (cTnT); High-Sensitivity Troponin (hs-cTnI, hs-cTnT); Cardiac Troponin; hs-Troponin

\section*{Principle}
Sandwich immunoassays use cardiac-specific antibodies to capture troponin I or T and generate a chemiluminescent or other instrument-specific signal. High-sensitivity classification depends on analytical imprecision and the proportion of healthy individuals whose concentrations can be measured above the assay's limit of detection; there is no single universal concentration or cutoff that applies to all hs-cTn assays. Results are normally reported in ng/L, with the assay-specific 99th-percentile upper reference limit and clinical decision thresholds supplied by the laboratory.

\section*{History}
Troponin was discovered as part of the contractile machinery of striated muscle in the 1960s. Cardiac-specific isoforms were identified in the 1980s and commercial immunoassays introduced in the 1990s. The rise of troponin as the preferred biomarker accelerated through the 2000s; the First Universal Definition of MI (2007) endorsed troponin superiority. High-sensitivity assays (Roche hs-cTnT, Abbott hs-cTnI) appeared in 2010s and became routine in the U.S. by 2017.

\section*{Why This Test Is Ordered}
\begin{itemize}
  \item Diagnosis of acute myocardial infarction in patients presenting with chest pain or ischemic equivalent symptoms
  \item Risk stratification and decision-making in acute coronary syndromes (STEMI, NSTEMI)
  \item Detection of myocardial injury in non-ACS conditions: sepsis, pulmonary embolism, takotsubo, severe heart failure, renal failure
  \item Monitoring of periprocedural myocardial injury after PCI, CABG, ablation
  \item Screening for cardiotoxicity in chemotherapy (anthracyclines, trastuzumab)
\end{itemize}

\section*{Is This a Primary or Secondary Test}
Primary biomarker for suspected MI; serial measurements required for diagnosis when baseline troponin is elevated or unclear.

\section*{How This Test Is Performed}
Venous blood is drawn into a laboratory-approved serum or plasma tube. For suspected ACS, samples are collected at presentation and repeated according to the validated local protocol, commonly at 1--2 hours with high-sensitivity assays or 3--6 hours with less sensitive assays. Stable patients with chronic elevations require interpretation of serial change, clinical context, and assay-specific decision limits rather than a single universal cutoff.

\section*{What Preparation Is Needed}
No fasting is required. Document symptom onset and the time of each sample; an early sample may be negative because the timing of release varies with the clinical event.

\section*{Contraindications and Precautions}
Troponin is cardiac-specific but not disease-specific: elevations occur in any condition that injures cardiomyocytes, including sepsis, renal failure, PE, stroke, stress cardiomyopathy, myocarditis, heart failure. Persistent chronic elevations in renal failure require interpretation with delta values. Heterophilic antibodies may falsely elevate troponin I (especially cTnI assays); confirmatory testing on different platforms if suspicion.

\section*{Risks}
Venipuncture only.

\section*{What Happens After the Test}
The result is reviewed with the ECG, symptoms, examination, imaging, and serial change. A rise and/or fall indicates acute myocardial injury; it does not distinguish type 1 MI, type 2 MI, or non-ischaemic injury by itself. Type 1 or type 2 MI requires evidence of ischemia in addition to the troponin pattern. Persistent elevation without a meaningful change may reflect chronic myocardial injury, including that associated with chronic kidney disease or structural heart disease.

\section*{Understanding Results}

\subsection*{Below Normal}
A low or undetectable troponin can support rule-out only when used within a validated assay-specific pathway and interpreted with symptom timing, ECG findings, and pre-test probability. A single negative result does not universally exclude MI.

\subsection*{Above Normal}
\begin{itemize}
  \item \textbf{Type 1 MI (plaque rupture, STEMI/NSTEMI):} Rise and/or fall of troponin with $\geq$ 1 value $>$ 99th percentile URL, plus ischemic symptoms, ECG changes, or imaging evidence
  \item \textbf{Type 2 MI (supply-demand mismatch):} Acute rise and/or fall in troponin above the assay-specific 99th percentile plus clinical evidence of ischemia caused by an imbalance in oxygen supply and demand (for example, tachyarrhythmia, hypotension, or severe anemia); plaque rupture is not required
  \item \textbf{Acute myocardial injury (non-MI)}: Sepsis, myocarditis, takotsubo, pulmonary embolism (right ventricular strain), acute heart failure
  \item \textbf{Chronic myocardial injury}: Persistently elevated troponin without delta; structural heart disease, renal failure, advanced heart failure
  \item \textbf{Chemotherapy cardiotoxicity}: Persistent elevation after anthracycline therapy
\end{itemize}

\section*{Normal Results}
\begin{itemize}
  \item \textbf{Reference limit:} The 99th-percentile upper reference limit is assay- and laboratory-specific; do not substitute a value from another platform.
  \item \textbf{Units:} High-sensitivity assays are generally reported in ng/L; conventional assays may use ng/mL. Unit conversion must be checked before interpretation.
  \item \textbf{Myocardial injury:} At least one troponin value above the assay-specific 99th-percentile upper reference limit.
  \item \textbf{MI:} Acute myocardial injury (rise and/or fall) plus clinical evidence of myocardial ischemia; troponin elevation alone is not sufficient.
  \item \textbf{ESC 0/1-h algorithm delta thresholds}: Based on absolute change at 1 hour (assay-specific)
\end{itemize}

\section*{Follow-up Tests}
\begin{itemize}
  \item \textbf{ECG (serial):} ST elevation, T-wave changes, bundle branch block
  \item \textbf{Echocardiography (urgent):} Regional wall motion abnormality in acute MI; right ventricular strain in PE; apical ballooning in takotsubo
  \item \textbf{ Serial troponin at 1, 2, and 3--6 hours}: Define delta pattern
  \item \textbf{Coronary angiography}: For STEMI or unstable NSTEMI patients
  \item \textbf{BNP/NT-proBNP}: Concurrent heart failure assessment
  \item \textbf{D-dimer, CT pulmonary angiography}: If PE suspected
  \item \textbf{CK-MB (historic; rarely needed)}: When troponin ambiguous or suspected reinfarction
\end{itemize}

\section*{References}
\begin{enumerate}
  \item Thygesen K, Alpert JS, Jaffe AS, et al. Fourth Universal Definition of Myocardial Infarction. J Am Coll Cardiol. 2018;72(18):2231-2264.
  \item Collet JP, Thiele H, Barbato E, et al. 2020 ESC guidelines for the management of acute coronary syndromes. Eur Heart J. 2021;42(14):1289-1367.
  \item Shah ASV, Anand A, Sandoval Y, et al. High-sensitivity cardiac troponin I at presentation. Lancet. 2015;386(10012):2481-2488.
  \item Morrow DA, et al. APACE-Troponin T risk score. Circulation. 2014;129(1):56-65.
\end{enumerate}

\clearpage
